\section{Un ejemplo de convergencia condicional}

\begin{eg}
  La serie \(\sum_{n \geq 1} \dfrac{(-1)^{n+1}}{n}\) converge, pero no lo hace absolutamente.
\end{eg}

\begin{proof}
  No converge absolutamente porque \(\sum \frac{1}{n}\) diverge. Veamos que converge.

  \emph{(i) Sumas parciales pares.} Se tiene
  \[
    S_{2n} = \left( 1 - \frac{1}{2} \right) + \left( \frac{1}{3} - \frac{1}{4} \right) + \cdots + \left( \frac{1}{2n-1} - \frac{1}{2n} \right),
  \]
  de modo que \(S_2 < S_4 < \cdots < S_{2n} < \cdots\) es creciente. Además,
  \[
    S_{2n} = 1 - \left( \frac{1}{2} - \frac{1}{3} \right) - \cdots - \left( \frac{1}{2n-2} - \frac{1}{2n-1} \right) - \frac{1}{2n} < 1,
  \]
  luego \({(S_{2n})}\) es creciente y acotada, y por lo tanto converge a algún \(\tilde{s}\).

  \emph{(ii) Sumas parciales impares.} Como \(S_{2n+1} = S_{2n} + \frac{1}{2n+1}\) y
  \[
    S_{2n+1} = 1 - \left( \frac{1}{2} - \frac{1}{3} \right) - \cdots - \left( \frac{1}{2n} - \frac{1}{2n+1} \right),
  \]
  se tiene \(S_1 > S_3 > \cdots > S_{2n+1} > \cdots > 0\), es decir, \({(S_{2n+1})}\) es decreciente y acotada, y converge a algún \(s'\).

  Como \(S_{2n+1} - S_{2n} = \frac{1}{2n+1} \to 0\), resulta \(\tilde{s} = s'\). Entonces \({(S_n)}\) converge, pues sus subsucesiones de índices pares e impares convergen al mismo límite.
\end{proof}

\section{Corolarios del criterio de comparación}

\begin{corollary}
  Sea \(\sum b_n\) una serie convergente de términos positivos. Si existen \(k > 0\) y \(n_0 \in \N\) tales que \(|a_n| \leq k\, b_n\) para todo \(n > n_0\), entonces \(\sum a_n\) converge absolutamente.
\end{corollary}

\begin{proof}
  Es una aplicación directa del criterio de comparación a \(|a_n|\) y \(b_n\).
\end{proof}

\begin{corollary}
  Si existen \(k > 0\), \(0 < c < 1\) y \(n_0 \in \N\) tales que \(|a_n| \leq k\, c^n\) para todo \(n > n_0\), entonces \(\sum a_n\) converge absolutamente.
\end{corollary}

\begin{proof}
  Como \(0 < c < 1\), la serie geométrica \(\sum c^n\) converge, y se aplica el corolario anterior.
\end{proof}

\begin{note}
  Tomando \(k = 1\), la condición \(|a_n| \leq c^n\) equivale a \(\sqrt[n]{|a_n|} \leq c < 1\). Si esto vale para todo \(n > n_0\), entonces \(\limsup_{n \to \infty} \sqrt[n]{|a_n|} \leq c < 1\). Recíprocamente, si \(\limsup_{n \to \infty} \sqrt[n]{|a_n|} < 1\), existen \(c < 1\) y \(n_0\) con \(\sqrt[n]{|a_n|} < c\) para \(n > n_0\).
\end{note}

\section{Criterios de convergencia}

\begin{theorem}[Criterio de la raíz]
  Sea \(\sum a_n\) una serie y \(\alpha = \limsup_{n \to \infty} \sqrt[n]{|a_n|}\).
  \begin{enumerate}[label=(\arabic*)]
    \item Si \(\alpha < 1\), entonces \(\sum a_n\) converge absolutamente.
    \item Si \(\alpha > 1\), entonces \(\sum a_n\) diverge.
    \item Si \(\alpha = 1\), el criterio no decide: tanto \(\sum \frac{1}{n}\) (que diverge) como \(\sum \frac{1}{n^2}\) (que converge) cumplen \(\sqrt[n]{|a_n|} \to 1\).
  \end{enumerate}
\end{theorem}

\begin{proof}
  (1) Sea \(\beta \in (\alpha, 1)\). Como \(\alpha = \inf_n \sup_{k \geq n} \sqrt[k]{|a_k|}\), existe \(N\) tal que \(\sqrt[n]{|a_n|} < \beta\) para todo \(n > N\), es decir, \(|a_n| < \beta^n\). Como \(0 < \beta < 1\), la serie \(\sum \beta^n\) converge y, por el criterio de comparación, \(\sum a_n\) converge absolutamente.

  (2) Si \(\alpha > 1\), existen infinitos índices con \(\sqrt[n]{|a_n|} > 1\) (si la sucesión de raíces no es acotada, también los hay), y entonces \(|a_n| > 1\) para infinitos \(n\). Luego \(a_n \not\to 0\) y la serie diverge.

  (3) Se tiene \(\sqrt[n]{1/n} \to 1\) por un ejemplo anterior y, por la misma razón, \(\sqrt[n]{1/n^2} = {(\sqrt[n]{1/n})}^2 \to 1\). La serie \(\sum \frac{1}{n}\) diverge, mientras que \(\sum \frac{1}{n^2}\) converge (por ejemplo, por el criterio de condensación, pues \(\sum 2^k \frac{1}{4^k} = \sum \frac{1}{2^k}\) converge).
\end{proof}

\begin{theorem}[Criterio de comparación de cocientes]
  Sea \(\sum a_n\) una serie con \(a_n \neq 0\) para todo \(n\), y sea \(\sum b_n\) una serie convergente de términos positivos. Si existe \(n_0 \in \N\) tal que
  \[
    \frac{|a_{n+1}|}{|a_n|} \leq \frac{b_{n+1}}{b_n} \quad \text{para todo } n > n_0,
  \]
  entonces \(\sum a_n\) converge absolutamente.
\end{theorem}

\begin{proof}
  Multiplicando las desigualdades correspondientes a \(n_0 + 1, n_0 + 2, \dots, n - 1\) se obtiene, para \(n > n_0 + 1\),
  \[
    \frac{|a_n|}{|a_{n_0 + 1}|} \leq \frac{b_n}{b_{n_0 + 1}}, \qquad \text{es decir,} \qquad |a_n| \leq \frac{|a_{n_0 + 1}|}{b_{n_0 + 1}} \, b_n.
  \]
  Por el criterio de comparación, \(\sum |a_n|\) converge.
\end{proof}

\begin{corollary}[Criterio de D'Alembert]
  Si existe \(c \in (0, 1)\) tal que \(\dfrac{|a_{n+1}|}{|a_n|} \leq c\) para todo \(n > n_0\), entonces \(\sum a_n\) converge absolutamente. En particular, esto ocurre si \(\limsup \dfrac{|a_{n+1}|}{|a_n|} < 1\). Si el cociente vale \(1\) el criterio no decide, y si \(\dfrac{|a_{n+1}|}{|a_n|} \geq 1\) para todo \(n > n_0\) la serie diverge (pues \(a_n \not\to 0\)).
\end{corollary}

\begin{proof}
  Se toma \(b_n = c^n\) en el teorema anterior, pues \(\sum c^n\) converge si \(c \in (0, 1)\). Si \(\limsup \frac{|a_{n+1}|}{|a_n|} < 1\), se elige \(c\) entre ese valor y \(1\) como en el criterio de la raíz.
\end{proof}

\begin{theorem}
  Sea \({(a_n)}\) una sucesión de términos positivos tal que \({\left( \frac{a_{n+1}}{a_n} \right)}\) es acotada. Entonces
  \[
    \liminf_{n \to \infty} \frac{a_{n+1}}{a_n} \leq \liminf_{n \to \infty} \sqrt[n]{a_n} \leq \limsup_{n \to \infty} \sqrt[n]{a_n} \leq \limsup_{n \to \infty} \frac{a_{n+1}}{a_n}.
  \]
  En particular, si existe \(\lim_{n \to \infty} \frac{a_{n+1}}{a_n}\), entonces existe \(\lim_{n \to \infty} \sqrt[n]{a_n}\) y ambos son iguales.
\end{theorem}

\begin{proof}
  La desigualdad del medio es evidente. Veamos que \(\limsup \sqrt[n]{a_n} \leq \limsup \frac{a_{n+1}}{a_n}\). Supongamos que no es así; entonces existe \(c\) con
  \[
    \limsup_{n \to \infty} \frac{a_{n+1}}{a_n} < c < \limsup_{n \to \infty} \sqrt[n]{a_n}.
  \]
  Por la definición de \(\limsup\), existe \(n_0\) tal que \(\frac{a_{n+1}}{a_n} < c\) para todo \(n \geq n_0\). Multiplicando estas desigualdades para \(n_0, \dots, n - 1\),
  \[
    \frac{a_n}{a_{n_0}} < c^{n - n_0} \quad \Longrightarrow \quad a_n < k\, c^n, \qquad k = \frac{a_{n_0}}{c^{n_0}}.
  \]
  Como \(\sqrt[n]{k} \to 1\), resulta
  \[
    \limsup_{n \to \infty} \sqrt[n]{a_n} \leq \limsup_{n \to \infty} \sqrt[n]{k\, c^n} = c \cdot \lim_{n \to \infty} \sqrt[n]{k} = c,
  \]
  lo cual es absurdo. La primera desigualdad se prueba de manera análoga, invirtiendo las desigualdades.

  Si el límite del cociente existe, \(\liminf\) y \(\limsup\) del cociente coinciden, y entonces los de la raíz también coinciden con él, de modo que la sucesión de raíces converge a ese valor.
\end{proof}

\begin{eg}
  Puede existir el límite de las raíces sin que exista el del cociente.
\end{eg}

\begin{proof}
  Sean \(0 < a < b\) y sea \({(x_n)}\) la sucesión que se obtiene multiplicando alternadamente por \(a\) y por \(b\):
  \[
    x_1 = a, \quad x_2 = a b, \quad x_3 = a^2 b, \quad x_4 = a^2 b^2, \quad \dots
  \]
  Entonces
  \[
    \frac{x_{n+1}}{x_n} =
    \begin{cases}
      b & \text{si } n \text{ es impar}, \\
      a & \text{si } n \text{ es par},
    \end{cases}
  \]
  de modo que no existe \(\lim_{n \to \infty} \frac{x_{n+1}}{x_n}\) (su \(\limsup\) es \(b\) y su \(\liminf\) es \(a\)). Por otro lado, \(x_{2k} = a^k b^k\) y \(x_{2k-1} = a^k b^{k-1}\), de modo que
  \[
    \sqrt[2k]{x_{2k}} = \sqrt{ab}, \qquad \sqrt[2k-1]{x_{2k-1}} = a^{\frac{k}{2k-1}} b^{\frac{k-1}{2k-1}} = \sqrt{ab} \cdot a^{\frac{1}{2(2k-1)}} b^{-\frac{1}{2(2k-1)}} \longrightarrow \sqrt{ab}.
  \]
  Como ambas subsucesiones convergen a \(\sqrt{ab}\), se tiene \(\lim_{n \to \infty} \sqrt[n]{x_n} = \sqrt{ab}\).
\end{proof}

\begin{theorem}[Criterio de Dirichlet]
  Sea \(\sum a_n\) una serie (no necesariamente convergente) cuyas sumas parciales \(S_n\) forman una sucesión acotada, y sea \({(b_n)}\) una sucesión decreciente de números positivos con \(\lim_{n \to \infty} b_n = 0\). Entonces \(\sum a_n b_n\) converge.
\end{theorem}

\begin{proof}
  Por sumación por partes, con \(S_0 = 0\),
  \begin{align*}
    a_1 b_1 + \cdots + a_n b_n & = S_1 (b_1 - b_2) + S_2 (b_2 - b_3) + \cdots + S_{n-1} (b_{n-1} - b_n) + S_n b_n \\
                               & = \sum_{i=1}^{n-1} S_i (b_i - b_{i+1}) + S_n b_n.
  \end{align*}
  Por hipótesis existe \(k > 0\) con \(|S_n| \leq k\) para todo \(n\). Como \(b_i - b_{i+1} \geq 0\), y la serie telescópica \(\sum_{i \geq 1} (b_i - b_{i+1})\) converge (a \(b_1\)), la serie \(\sum S_i (b_i - b_{i+1})\) converge absolutamente por comparación, y en particular converge. Además \(S_n b_n \to 0\), por ser \((S_n)\) acotada y \(b_n \to 0\). Luego la sucesión de sumas parciales de \(\sum a_n b_n\) converge.
\end{proof}

\begin{corollary}[Criterio de Abel]
  Si \(\sum a_n\) converge y \({(b_n)}\) es una sucesión decreciente de términos positivos, entonces \(\sum a_n b_n\) converge.
\end{corollary}

\begin{proof}
  Como \({(b_n)}\) es decreciente y acotada inferiormente por \(0\), converge a algún \(c \geq 0\). La sucesión \({(b_n - c)}\) es decreciente, de términos no negativos, y tiende a \(0\). Como \(\sum a_n\) converge, sus sumas parciales son acotadas, y por el criterio de Dirichlet \(\sum a_n (b_n - c)\) converge. Entonces
  \[
    \sum a_n b_n = \sum a_n (b_n - c) + c \sum a_n
  \]
  es suma de dos series convergentes.
\end{proof}

\begin{corollary}[Criterio de Leibniz]
  Si \({(b_n)}\) es decreciente y \(b_n \to 0\), entonces \(\sum (-1)^n b_n\) converge.
\end{corollary}

\begin{proof}
  La serie \(\sum (-1)^n\) no converge, pero sus sumas parciales toman los valores \(-1\) y \(0\), y son acotadas. Se aplica el criterio de Dirichlet.
\end{proof}

\section{Parte positiva y parte negativa}

\begin{definition}
  Dada una serie \(\sum a_n\), para cada \(n \in \N\) definimos
  \[
    p_n =
    \begin{cases}
      a_n & \text{si } a_n > 0,    \\
      0   & \text{en otro caso},
    \end{cases}
    \qquad
    q_n =
    \begin{cases}
      -a_n & \text{si } a_n < 0,    \\
      0    & \text{en otro caso}.
    \end{cases}
  \]
  Entonces \(p_n, q_n \geq 0\), \(a_n = p_n - q_n\) y \(|a_n| = p_n + q_n = a_n + 2 q_n = 2 p_n - a_n\).
\end{definition}

\begin{note}
  Si \(\sum a_n\) converge absolutamente, para todo \(k\) se tiene
  \[
    \sum_{n \geq 1} |a_n| \geq \sum_{n=1}^{k} |a_n| = \sum_{n=1}^{k} p_n + \sum_{n=1}^{k} q_n,
  \]
  de modo que las sumas parciales de \(\sum p_n\) y \(\sum q_n\) son crecientes y acotadas por \(\sum |a_n|\), y ambas series convergen. Recíprocamente, si \(\sum p_n\) y \(\sum q_n\) convergen, entonces \(\sum |a_n| = \sum (p_n + q_n)\) converge, y \(\sum a_n\) converge absolutamente.
\end{note}

\begin{theorem}
  Si \(\sum a_n\) converge condicionalmente, entonces \(\sum p_n\) y \(\sum q_n\) divergen.
\end{theorem}

\begin{proof}
  Supongamos que \(\sum q_n\) converge, con suma \(c\). Como \(|a_n| = a_n + 2 q_n\),
  \[
    \sum_{n=1}^{k} |a_n| = \sum_{n=1}^{k} a_n + 2 \sum_{n=1}^{k} q_n \longrightarrow \sum_{n \geq 1} a_n + 2c \quad (k \to \infty),
  \]
  de modo que \(\sum |a_n|\) converge, lo cual es absurdo pues \(\sum a_n\) no converge absolutamente. Luego \(\sum q_n\) diverge. Análogamente, usando \(|a_n| = 2 p_n - a_n\), se prueba que \(\sum p_n\) diverge.
\end{proof}

\section{Reordenamientos}

\begin{definition}
  Reordenar una serie \(\sum a_n\) significa tomar una biyección \(\varphi\colon \N \to \N\) y considerar la serie \(\sum b_n\) con \(b_n = a_{\varphi(n)}\).
\end{definition}

\begin{theorem}
  Todo reordenamiento de una serie absolutamente convergente converge al mismo valor que la serie original.
\end{theorem}

\begin{proof}
  \emph{Caso \(a_n \geq 0\).} Sea \(\varphi\colon \N \to \N\) una biyección y \(b_n = a_{\varphi(n)}\). Sean \(S_n\) y \(T_n\) las sumas parciales de \(\sum a_n\) y \(\sum b_n\). Para cada \(n\), sea \(m = \max\{\varphi(1), \dots, \varphi(n)\}\), de modo que \(\{\varphi(1), \dots, \varphi(n)\} \subset \{1, \dots, m\}\). Como \(a_j \geq 0\),
  \[
    T_n = \sum_{i=1}^{n} a_{\varphi(i)} \leq \sum_{j=1}^{m} a_j = S_m \leq \sum_{j \geq 1} a_j.
  \]
  Entonces \(\sum b_n\) converge y \(\sum b_n \leq \sum a_n\). Como \(a_n = b_{\varphi^{-1}(n)}\) es un reordenamiento de \(\sum b_n\), el mismo argumento da \(\sum a_n \leq \sum b_n\). Luego son iguales.

  \emph{Caso general.} Como \(\sum a_n\) converge absolutamente, \(\sum a_n = \sum p_n - \sum q_n\) con \(\sum p_n\) y \(\sum q_n\) convergentes y de términos no negativos. Las partes positiva y negativa de \(b_n = a_{\varphi(n)}\) son \(p_{\varphi(n)}\) y \(q_{\varphi(n)}\), que son reordenamientos de \(\sum p_n\) y \(\sum q_n\). Por el caso anterior,
  \[
    \sum b_n = \sum p_{\varphi(n)} - \sum q_{\varphi(n)} = \sum p_n - \sum q_n = \sum a_n. \qedhere
  \]
\end{proof}

\begin{theorem}[Riemann]
  Sea \(\sum a_n\) una serie condicionalmente convergente y sea \(c \in \R\). Entonces existe un reordenamiento \(\sum b_n\) de \(\sum a_n\) tal que \(\sum b_n = c\).
\end{theorem}

\begin{proof}
  Como \(\sum a_n\) converge, \(a_n \to 0\). Sean \(P_1, P_2, \dots\) los términos positivos de la serie y \(-Q_1, -Q_2, \dots\) los negativos, cada uno en su orden original. Entonces \(P_k \to 0\) y \(Q_k \to 0\), y por el teorema anterior \(\sum P_k = \sum Q_k = +\infty\). (Los términos nulos, si los hay, pueden intercalarse en cualquier lugar sin alterar el argumento.)

  Construimos el reordenamiento por etapas:
  \begin{enumerate}[label=(\arabic*)]
    \item Se toman \(P_1, P_2, \dots\) hasta que la suma parcial supere por primera vez a \(c\).
    \item Se agregan \(-Q_1, -Q_2, \dots\) hasta que la suma parcial sea por primera vez menor que \(c\).
    \item Se agregan los siguientes \(P_k\) hasta superar nuevamente a \(c\), luego los siguientes \(-Q_k\) hasta quedar por debajo de \(c\), y así sucesivamente.
  \end{enumerate}
  El proceso nunca se detiene, pues \(\sum P_k\) y \(\sum Q_k\) divergen, y usa todos los términos de la serie original, ya que en cada etapa se consumen términos nuevos de cada tipo.

  Veamos que las sumas parciales \(T_N\) convergen a \(c\). Al terminar una etapa de suma de positivos, con último término \(P_j\), la suma parcial supera a \(c\) pero antes de sumar \(P_j\) no lo hacía, de modo que \(0 < T - c \leq P_j\). Análogamente, al terminar una etapa de negativos con último término \(-Q_j\), \(0 < c - T \leq Q_j\). Dentro de cada etapa, las sumas parciales varían monótonamente entre dos de estos valores, de modo que están entre ambos. Entonces, para \(N\) en una etapa avanzada, \(|T_N - c| \leq \max\{P_j, Q_{j'}\}\), donde \(P_j\) y \(Q_{j'}\) son los últimos términos usados en los extremos de dicha etapa. Como \(P_j \to 0\) y \(Q_{j'} \to 0\) cuando el número de etapas tiende a infinito, resulta \(T_N \to c\).
\end{proof}
